% =============================================================================
%   Discrete adjoints for RBF-FD shape optimization in Macchiato.jl
%   Theory and implementation of the manual gradient pipeline
%
%   Compile with:  pdflatex report_phase_D_L2_helmholtz   (twice for refs)
%   Or:            latexmk -pdf report_phase_D_L2_helmholtz.tex
%   Or paste into Overleaf — all packages are standard.
% =============================================================================
\documentclass[11pt,a4paper]{article}

\usepackage[margin=1in]{geometry}
\usepackage{amsmath,amssymb,amsthm,bm}
\usepackage{graphicx}
\graphicspath{{docs/}{figures/}}    % look here when compiling from the repo root
\usepackage{booktabs}
\usepackage{array}
\usepackage{caption}
\usepackage{subcaption}
\usepackage{hyperref}
\usepackage{xcolor}
\usepackage{listings}
\usepackage{algorithm}
\usepackage{algpseudocode}
\usepackage{enumitem}

\hypersetup{
  colorlinks=true,
  linkcolor=blue!50!black,
  citecolor=blue!50!black,
  urlcolor=blue!50!black,
}

\lstdefinelanguage{Julia}{
  morekeywords={function,end,if,else,elseif,for,while,return,break,continue,
    struct,mutable,abstract,type,module,using,import,export,const,let,do,
    begin,finally,try,catch,throw,global,local,macro,quote,
    true,false,nothing,missing,in,isa,where,as,is,outer,inline},
  sensitive=true,
  morecomment=[l]{\#},
  morecomment=[n]{\#=}{=\#},
  morestring=[b]",
  morestring=[b]',
  morestring=[s]{"""}{"""},
}

\lstset{
  basicstyle=\ttfamily\small,
  breaklines=true,
  columns=fullflexible,
  keepspaces=true,
  language=Julia,
  commentstyle=\color{gray},
  numbers=left,
  numberstyle=\tiny\color{gray},
  numbersep=8pt,
  frame=single,
  xleftmargin=2em,
}

\newcommand{\mat}[1]{\mathbf{#1}}
\newcommand{\vect}[1]{\boldsymbol{#1}}
\newcommand{\pts}{\vect{p}}
\newcommand{\norm}[1]{\left\lVert#1\right\rVert}
\newcommand{\inner}[2]{\langle #1,\,#2\rangle}
\newcommand{\reals}{\mathbb{R}}
\newcommand{\dd}{\,\mathrm{d}}
\newcommand{\D}{\mathcal{D}}
\newcommand{\Loss}{\mathcal{L}}
\newcommand{\TODO}[1]{\textcolor{red}{[TODO: #1]}}

\title{
  \textbf{Discrete adjoints for RBF-FD shape optimization in
  \texttt{Macchiato.jl}}\\[0.6em]
  \large Theory and implementation of the manual gradient pipeline
}
\author{Davide Miotti}
\date{2026-06-02 (working document)}

\begin{document}
\maketitle

\begin{abstract}
We document the full discrete adjoint pipeline implemented in
\texttt{Macchiato.jl} for shape optimization on RBF-FD meshless
discretizations of elasticity. Starting from a continuous
shape-optimization problem on a parameterized domain, we develop a
five-step manual adjoint that decomposes the sensitivity calculation
into a forward elasticity solve caching RBF stencils, a single adjoint
solve, extraction of sensitivities at the level of weight-matrix
entries, a pure-Julia backward pass through the RBF stencil
computations, and the assembly of explicit external contributions. Each
step has closed-form analytic content and is validated to machine
precision against fifth-order central finite differences. We give
detailed derivations for each ingredient: the interior elasticity
rows for plane stress, the Neumann row family with traction BCs, the
shape-dependent dead-load contribution (Phase D Level 1), the
differentiable polyline normal model (Phase D Level 2), and the
Helmholtz-PDE sensitivity filter on the boundary loop. A 2D cantilever
benchmark illustrates the pipeline end-to-end. A plate-with-hole study
then motivates reorganizing the optimizer around a Fourier design
parameterization and an actively re-meshed cloud: the boundary gradient
noise is traced to the geometry sensitivity of the RBF-FD weights
($\partial\mat W/\partial\pts$), frequency filtering is shown unable to
cure it, and a clean-DOF design space with a two-front (morph plus
re-anchored re-mesh) cloud controller takes its place. We close by mapping the
open extensions --- generalization beyond linear elasticity,
integration with the geometry kernel in
\texttt{WhatsThePoint.jl}, scale-up, and the path to the longer-term
multi-physics framework --- against established state-of-the-art
positioning.
\end{abstract}

\paragraph{Status of this report.}
This is a working technical document, written to support both
ongoing development and a future engineering-oriented publication.
The cantilever example used for illustration is under active
development at the time of writing: in particular, an artefact in
the tip-shrinkage gradient was found to have been masked by a
previous corner-freeze convention, and is being repaired. The
\emph{numerical} compliance figures cited in
Section~\ref{sec:cantilever} should therefore be regarded as
illustrative snapshots; the qualitative behaviour they convey
(regularization effect, role of differentiable normals) is robust
and is what we use them to demonstrate. Note also that the gradient
filter of Section~\ref{sec:filter}, presented there as the
regularization cure, is \emph{superseded} by Section~\ref{sec:plate}:
the plate-with-hole study shows frequency filtering cannot separate the
broadband weight-sensitivity artifact from the signal, and the cure is
the design-space-plus-re-meshing architecture developed there.

\tableofcontents
\bigskip

\section{Introduction}
\label{sec:intro}

Shape optimization treats the shape of a physical domain as a design
variable and asks for the geometry that extremizes a physical
objective subject to physical constraints. Computing the gradient of
the objective with respect to the design — the \emph{sensitivity} —
is the central computational task, and the dominant strategy for
problems with many design degrees of freedom is the \emph{adjoint
method}: a single auxiliary linear solve (the ``adjoint'') replaces
many directional forward solves that direct differentiation would
require.

\texttt{Macchiato.jl} implements this strategy for radial-basis-function
finite-difference (RBF-FD) discretizations. RBF-FD is a meshless
method: the discretization is a point cloud rather than a connected
mesh, and the differential operators (Laplacian, divergence, mixed
partial, etc.) are realized as sparse weight matrices built by local
RBF stencil fitting. The meshless choice has structural advantages for
shape optimization — discretization smoothly tracks geometry
deformation, mesh tangling is not a concept, topology change is
absorbed naturally — but it puts a premium on \emph{differentiability
of the discretization itself}: the weight matrices, the boundary
normals, the per-row coefficient blocks must all be differentiable
functions of the point positions, and the adjoint must thread the
gradient cleanly through every dependence.

The pipeline reported here is a discrete adjoint whose chain rule is
\emph{orchestrated by hand} --- but this must be stated precisely, because the
differentiation of the RBF basis is \emph{not} done by hand. The RBF-FD weight
matrices are built by \texttt{RadialBasisFunctions.jl}, which differentiates the
basis functions automatically to apply each differential operator; we never write a
basis derivative ourselves. What is manual is the \emph{shape adjoint}: the chain
rule that carries the loss sensitivity back through the assembled discrete system to
the node coordinates --- the five-step pull-back of Section~\ref{sec:five-steps} ---
is written explicitly rather than obtained by tracing a reverse-mode AD engine
through the whole forward solve. The motivation is twofold. First, AD-tracing the
full solve (the UMFPACK factorization and the local RBF saddle solves) incurs heavy
LLVM compile cost: an earlier Mooncake attempt was abandoned at five-minute compiles
even at $N=25$ points (Section~\ref{sec:diffstrategy} discusses how custom pull-backs
would recover a clean AD-native API without that cost). Second, the explicit form
makes each component inspectable, FD-validatable to machine precision, and reusable
across physics: the same Step~2 adjoint solve, Step~4 stencil pull-back, and Step~3
coefficient-extraction pattern serve every PDE.

This document presents the full chain. Section~\ref{sec:problem} fixes
notation. Section~\ref{sec:discretization} specifies the RBF-FD
discretization of plane-stress linear elasticity, including the
interior assembly, the Dirichlet rows, and the four-coefficient
Neumann row family. Section~\ref{sec:adjoint-theory} derives the
discrete adjoint from first principles, exposing the
$\partial\mat A/\partial\pts$ and $\partial\vect b/\partial\pts$
decompositions that drive every later step.
Section~\ref{sec:five-steps} walks the five-step pipeline in detail,
with the elasticity coefficient extraction in Step~3 worked out
explicitly. Section~\ref{sec:boundary-diff} treats the differentiable
Neumann data (Phase D Levels 1 and 2): shape-dependent traction
values and shape-dependent boundary normals.
Section~\ref{sec:filter} treats sensitivity filtering on the boundary
loop, both the simple 3-point average and the Helmholtz-PDE filter.
Section~\ref{sec:verification} describes the verification methodology.
Section~\ref{sec:cantilever} runs the pipeline on a 2D cantilever
example. Section~\ref{sec:plate} reorganizes the optimizer around a Fourier
design space and active re-meshing for the plate-with-hole problem.
Section~\ref{sec:diffstrategy} discusses the differentiation strategy: the
manual analytic pull-backs and the path to an AD-native \texttt{gradient} API.
Section~\ref{sec:generalization} maps the extensions ---
to other elliptic PDEs, to nonlinear systems, to multi-physics, and
to 3D. Section~\ref{sec:sota} positions the work relative to the
state of the art. Section~\ref{sec:openwork} lists open work.

\section{The method in outline}
\label{sec:outline}

Before specializing to elasticity, we state the approach in PDE-agnostic terms.
Shape optimization here is a loop over a \emph{smooth, low-dimensional design}
$\vect l$ (Fourier or surface-harmonic modes, or CAD / RBF control points); the PDE
is solved on a point cloud that discretizes the design's geometry, and the gradient
is taken by a discrete adjoint and then \emph{contracted} onto the design.

\begin{footnotesize}
\begin{verbatim}
  l --geometry--> pts --RBF-FD assemble--> (A,b) --solve--> u --loss--> L
  ^                                                                     |
  |                                                                     v
  +-- step <-- contract (dpts/dl)^T <-- discrete adjoint: dL/dpts <-----+
                                        (one forward + one adjoint solve)

  the cloud is kept healthy by TWO FRONTS:
    Front 1 (within a step):  morph the interior to follow the boundary
                              (smooth, differentiable; adjoint via transpose)
    Front 2 (between steps):  remesh + re-anchor when a quality indicator trips
\end{verbatim}
\end{footnotesize}

The four ingredients, and where each is treated:
\begin{enumerate}[label=(\arabic*)]
  \item \emph{Design space} $\vect l \mapsto \pts$ --- a smooth parameterization
        whose contraction~\eqref{eq:design-contraction} removes the mesh-scale
        gradient noise (Section~\ref{sec:plate}). PDE-independent.
  \item \emph{Forward solve} --- the PDE discretized by RBF-FD on the cloud
        (Section~\ref{sec:discretization}, for elasticity).
  \item \emph{Discrete adjoint} $\partial\Loss/\partial\pts$ --- one adjoint solve
        plus an analytic back-projection through the RBF stencils
        (Sections~\ref{sec:adjoint-theory}--\ref{sec:five-steps}). This is the only
        PDE-specific machinery; everything else is reused across physics.
  \item \emph{Two-front cloud control} --- morph and re-anchored remesh keep the
        discretization healthy and the discrete optimum unbiased
        (Section~\ref{sec:plate-twofront}). PDE-independent.
\end{enumerate}

Sections~\ref{sec:problem}--\ref{sec:boundary-diff} fill in ingredients (2)--(3) for
linear elasticity; Section~\ref{sec:plate} demonstrates (1) and (4) on the
plate-with-hole problem; Section~\ref{sec:diffstrategy} discusses how the whole
chain could be exposed behind a single \texttt{gradient} call.

\section{Continuous problem statement}
\label{sec:problem}

\subsection{Domain and design parameterization}

Let $\Omega(\vect l) \subset \reals^d$ ($d \in \{2,3\}$) be a bounded
domain whose shape is determined by a design vector
$\vect l \in \reals^{n_l}$. The boundary
$\partial\Omega = \Gamma_D \cup \Gamma_N$ partitions into disjoint
parts where Dirichlet and Neumann data are prescribed. The shape
dependence enters through any of three layers, which we keep separate
throughout this document because they enter the gradient through
distinct paths:

\begin{enumerate}[label=(\roman*)]
  \item \emph{Boundary-node positions} $\pts_b(\vect l)$ --- the
        discrete-level identity of the boundary.
  \item \emph{Interior-node positions} $\pts_i(\vect l)$ --- generated
        from $\pts_b$ by a node-distribution policy (see
        Section~\ref{sec:openwork}).
  \item \emph{Boundary normals and traction data}
        $\vect n(\vect l)$, $\vect t(\vect l)$ --- explicit functions
        of the boundary geometry whenever the loading is
        shape-coupled (e.g.\ parabolic shear on a free end whose
        length is a design parameter, or follower pressure).
\end{enumerate}

It is essential to keep two layers distinct, because conflating them is the
central pitfall of meshless shape optimization (diagnosed in
Section~\ref{sec:plate}): the \emph{design} is the low-dimensional, smooth shape
parameterization $\vect l$ --- a few Fourier or surface-harmonic modes, or CAD /
RBF control points --- while the point cloud $\pts(\vect l) \in \reals^{dN}$ is
merely its \emph{discretization}. The adjoint developed in
Sections~\ref{sec:adjoint-theory}--\ref{sec:five-steps} delivers the
discretization-level sensitivity $\partial\Loss/\partial\pts$; this is an
\emph{intermediate} quantity. The usable gradient is its contraction onto the
design,
\begin{equation}
  \frac{\partial\Loss}{\partial\vect l}
   = \Big(\frac{\partial\pts}{\partial\vect l}\Big)^{\!\top}
     \frac{\partial\Loss}{\partial\pts},
  \label{eq:design-contraction}
\end{equation}
where $\partial\pts/\partial\vect l$ is the (cheap) parameterization Jacobian.
\textbf{Treating the node positions themselves as the design --- descending
directly on $\partial\Loss/\partial\pts$ --- fails}: that gradient is dominated by
a mesh-scale noise mode (Section~\ref{sec:plate-noise}), and the contraction
\eqref{eq:design-contraction} onto a smooth low-dimensional $\vect l$ is precisely
what removes it. We therefore carry $\partial\Loss/\partial\pts$ through the body of
the report as the \emph{core primitive} --- it is PDE-specific and is where all the
analytic work lives --- with the understanding that the design-space contraction is
applied on top and is what makes the optimizer converge. The geometry map
$\vect l \mapsto \pts$ is provided by the kernel in \texttt{WhatsThePoint.jl} and is
kept outside the PDE side of the gradient.

\subsection{Linear elasticity}

The strong form, for plane-stress isotropic linear elasticity in 2D
with Lamé parameters $(\lambda^\star, \mu)$:
\begin{align}
  \nabla \!\cdot\! \vect\sigma(\vect u) \;&=\; \vect 0
    \quad &&\text{in } \Omega, \\
  \vect u \;&=\; \vect u_D
    \quad &&\text{on } \Gamma_D, \\
  \vect\sigma(\vect u)\!\cdot\!\vect n \;&=\; \vect t
    \quad &&\text{on } \Gamma_N,
\end{align}
with the constitutive law
$\vect\sigma(\vect u) = \lambda^\star(\nabla\!\cdot\!\vect u)\,\mat I
  + \mu(\nabla\vect u + \nabla\vect u^\top)$ and outward unit normal
$\vect n$. The displacement $\vect u: \Omega \to \reals^2$ is the
state.

\subsection{The shape-optimization problem}

Given an objective functional $\Loss(\vect u, \pts)$ and an optional
geometric constraint $V(\pts) = V_0$ (typically area or volume), the
shape-optimization problem is
\begin{equation}
  \min_{\pts}\; \Loss\big(\vect u(\pts),\, \pts\big)
  \quad \text{s.t.}\quad
  \begin{cases}
    \nabla\!\cdot\!\vect\sigma(\vect u) = \vect 0
      & \text{in } \Omega(\pts), \\
    \vect u = \vect u_D & \text{on } \Gamma_D, \\
    \vect\sigma\!\cdot\!\vect n = \vect t & \text{on } \Gamma_N, \\
    V(\pts) = V_0.
  \end{cases}
  \label{eq:shape-opt}
\end{equation}
For the illustrative problem of Section~\ref{sec:cantilever} we use the
compliance loss
$\Loss(\vect u, \pts) = \vect b(\pts)^\top \vect u$ with the area
constraint applied as a two-sided quadratic penalty
$\rho_{\text{pen}}\big(V(\pts) - V_0\big)^2$ added to the objective.
The treatment of the constraint is otherwise orthogonal to the
adjoint machinery developed here.

\section{RBF-FD discretization of elasticity}
\label{sec:discretization}

\subsection{Stencils and weight matrices}

We sample the domain with $N$ collocation points
$\{\pts_i\}_{i=1}^N \subset \overline\Omega$ and, for each one, build
a local neighborhood (``stencil'') of the $k$ nearest neighbors.
On each stencil we fit a polyharmonic-spline radial basis
$\phi(r) = r^m$ (we use $m=3$) augmented by polynomials up to degree
$p$ (we use $p=3$), and we solve a small saddle-point system that
returns the coefficients of any linear differential operator
$\mathcal L$ acting on the stencil basis. The resulting weights form
a row of the sparse \emph{weight matrix}
$\mat W^{\mathcal L} \in \reals^{N \times N}$: for each operator we
need (second partials $\partial_{xx}, \partial_{yy}, \partial_{xy}$
for elasticity interior assembly; first partials $\partial_x,
\partial_y$ for Neumann traction assembly), the same stencil layout
produces a different set of weights, but the sparsity pattern is
shared across operators on the same stencil.

The action $\mat W^{\mathcal L}$ on a discrete field
$\vect v \in \reals^N$ approximates $\mathcal L v$ at each sample
point:
\begin{equation}
  (\mat W^{\mathcal L} \vect v)_i
  \;\approx\;
  (\mathcal L v)(\pts_i),
  \qquad i = 1,\ldots,N.
\end{equation}
Two facts will be central to the adjoint construction: \emph{(i)}
$\mat W^{\mathcal L}$ is a smooth function of all $\pts_j$ in the
stencil of row $i$, hence of $\pts$ globally; \emph{(ii)} the
backward pass $\partial\mat W^{\mathcal L}/\partial\pts$ for given
gradient $\Delta\mat W^{\mathcal L}$ is implemented as a pure-Julia
routine \texttt{\_pullback\_weights!} in
\texttt{RadialBasisFunctions.jl} that reuses the saddle-point
factorization cached during the forward weight build, costing only
$O(k^3)$ per row.

\subsection{Block assembly of plane-stress elasticity}

The displacement state is interleaved as
$\vect u = (u_1,\ldots,u_N,\,v_1,\ldots,v_N)^\top \in \reals^{2N}$.
The Navier equations $\nabla\!\cdot\!\vect\sigma = \vect 0$ in plane
stress become, after expanding the constitutive law,
\begin{align}
  c_1 \partial_{xx} u + c_2 \partial_{yy} u + c_3 \partial_{xy} v
    \;&=\; 0, \\
  c_3 \partial_{xy} u + c_2 \partial_{xx} v + c_1 \partial_{yy} v
    \;&=\; 0,
\end{align}
with $c_1 = \lambda^\star + 2\mu$, $c_2 = \mu$, $c_3 = \lambda^\star + \mu$.
Substituting the discrete operators gives the $2N \times 2N$ block
matrix
\begin{equation}
  \mat A_{\text{int}}
  = \begin{bmatrix}
    c_1 \mat W^{\partial_{xx}} + c_2 \mat W^{\partial_{yy}}
    & c_3 \mat W^{\partial_{xy}} \\[2pt]
    c_3 \mat W^{\partial_{xy}}
    & c_2 \mat W^{\partial_{xx}} + c_1 \mat W^{\partial_{yy}}
  \end{bmatrix}.
  \label{eq:A-int}
\end{equation}
This block is applied at all interior rows. Boundary rows
(Dirichlet, Neumann) overwrite the corresponding rows of
$\mat A_{\text{int}}$ after assembly.

\subsection{Dirichlet rows}

For each Dirichlet point with global index $i \in \Gamma_D$, the
$x$-equation row $i$ and the $y$-equation row $i+N$ are overwritten:
\begin{equation}
  \mat A[i,:] = \vect e_i^\top, \quad
  \mat A[i{+}N,:] = \vect e_{i+N}^\top, \quad
  b_i = u_{D,i}^x,\quad b_{i+N} = u_{D,i}^y.
  \label{eq:dirichlet}
\end{equation}
These rows do not depend on $\pts$ at all, so they contribute zero
to $\partial\mat A/\partial\pts$ and $\partial\vect b/\partial\pts$.

\subsection{Neumann (traction) rows}
\label{sec:neumann-row}

For each Neumann point with local index
$i_{\text{loc}} \in \{1,\ldots,n_{\text{neu}}\}$ and global index
$g = \text{neumann\_ids}[i_{\text{loc}}]$, with outward unit normal
$\vect n = (n_x, n_y)$, stencil $\mathcal S$, and prescribed traction
$\vect t = (t_x, t_y)$, the discrete equations
$\vect\sigma\!\cdot\!\vect n = \vect t$ become, after expanding the
constitutive law and applying $\mat W^{\partial_x},
\mat W^{\partial_y}$ on the stencil:
\begin{align}
  \mat A[g,\,j]
    \;&=\; n_x(\lambda^\star{+}2\mu) \mat W^{\partial_x}[i_{\text{loc}},j]
        + n_y\mu\, \mat W^{\partial_y}[i_{\text{loc}},j], \notag\\
  \mat A[g,\,j{+}N]
    \;&=\; n_y\mu\, \mat W^{\partial_x}[i_{\text{loc}},j]
        + n_x \lambda^\star \mat W^{\partial_y}[i_{\text{loc}},j], \notag\\
  \mat A[g{+}N,\,j]
    \;&=\; n_y \lambda^\star \mat W^{\partial_x}[i_{\text{loc}},j]
        + n_x\mu\, \mat W^{\partial_y}[i_{\text{loc}},j], \notag\\
  \mat A[g{+}N,\,j{+}N]
    \;&=\; n_x\mu\, \mat W^{\partial_x}[i_{\text{loc}},j]
        + n_y(\lambda^\star{+}2\mu) \mat W^{\partial_y}[i_{\text{loc}},j],
  \label{eq:neumann-row}
\end{align}
for each $j \in \mathcal S$, with
$b_g = t_x$ and $b_{g+N} = t_y$. The pattern is: each Neumann row
entry is a linear combination of one or more weight-matrix entries,
with coefficients that depend on $(n_x, n_y, \lambda^\star, \mu)$.
We package this in a \texttt{TractionLayout} struct, storing for
each entry the pair $(c_x, c_y)$ that multiplies
$\big(\mat W^{\partial_x}[i_{\text{loc}}, j],
\mat W^{\partial_y}[i_{\text{loc}}, j]\big)$.

\subsection{The full system}

After applying~\eqref{eq:A-int}, \eqref{eq:dirichlet},
and~\eqref{eq:neumann-row}, the assembled system is
\begin{equation}
  \mat A(\pts)\, \vect u \;=\; \vect b(\pts),
  \label{eq:system}
\end{equation}
solved as a single LU factorization for the forward pass. Both
$\mat A$ and $\vect b$ depend on $\pts$:
\begin{itemize}[leftmargin=2em]
  \item $\mat A$ via the interior weights
        $\mat W^{\partial_{xx}}, \mat W^{\partial_{yy}},
        \mat W^{\partial_{xy}}$ (always);
  \item $\mat A$ via the Neumann weights
        $\mat W^{\partial_x}, \mat W^{\partial_y}$ (whenever
        $\Gamma_N$ is non-empty);
  \item $\mat A$ via the boundary normals $\vect n$ (whenever the
        normals are computed from the geometry rather than
        prescribed --- Phase D Level 2);
  \item $\vect b$ via the traction values $\vect t$ (whenever the
        loading is shape-coupled --- Phase D Level 1);
  \item $\vect b$ via interior body forces (none in the elasticity
        examples here, but the framework is set up for it).
\end{itemize}
The Dirichlet rows of $\mat A$ and $\vect b$ are independent of
$\pts$ by construction.

\section{The discrete adjoint: derivation}
\label{sec:adjoint-theory}

\subsection{Total derivative of the loss}

The loss has the structure
\begin{equation}
  \Loss(\vect u, \pts)
  = \Loss_{\text{state}}(\vect u, \pts) + \Loss_{\text{geom}}(\pts),
\end{equation}
where $\Loss_{\text{state}}$ depends on the state $\vect u$ and
optionally on $\pts$ directly, and $\Loss_{\text{geom}}$ depends on
$\pts$ alone (e.g.\ the area penalty). The total derivative with
respect to $\pts$ is
\begin{equation}
  \frac{\dd\Loss}{\dd\pts}
  \;=\;
  \underbrace{\frac{\partial\Loss}{\partial\vect u}\,
    \frac{\dd\vect u}{\dd\pts}}_{\text{state-mediated}}
  \;+\;
  \underbrace{\frac{\partial\Loss}{\partial\pts}}_{\text{explicit}}.
  \label{eq:total-deriv}
\end{equation}
The first term is what the adjoint method computes efficiently. The
second is assembled explicitly in Step~5
(Section~\ref{sec:step5-explicit}).

\subsection{Sensitivity of the state via the implicit function theorem}

Differentiating the linear system~\eqref{eq:system} with respect to
$\pts$:
\begin{equation}
  \frac{\partial\mat A}{\partial\pts}\,\vect u
  + \mat A \frac{\dd\vect u}{\dd\pts}
  \;=\; \frac{\partial\vect b}{\partial\pts},
\end{equation}
so
\begin{equation}
  \frac{\dd\vect u}{\dd\pts}
  \;=\;
  \mat A^{-1}\Big(\frac{\partial\vect b}{\partial\pts}
                  - \frac{\partial\mat A}{\partial\pts}\,\vect u\Big).
  \label{eq:du-dp}
\end{equation}
Substituting back into~\eqref{eq:total-deriv}:
\begin{equation}
  \frac{\dd\Loss}{\dd\pts}
  \;=\;
  \frac{\partial\Loss}{\partial\vect u}\, \mat A^{-1}
    \Big(\frac{\partial\vect b}{\partial\pts}
        - \frac{\partial\mat A}{\partial\pts}\,\vect u\Big)
  + \frac{\partial\Loss}{\partial\pts}.
\end{equation}

\subsection{The adjoint variable}

Naively this requires one linear solve per coordinate ---
$O(d N)$ solves --- because
$\partial\vect b/\partial\pts$ and $\partial\mat A/\partial\pts$
have a column per coordinate. The adjoint trick is to take the
left contraction first. Define
\begin{equation}
  \boxed{\;\vect\eta
    \;\equiv\; \mat A^{-\top}\,
      \left(\tfrac{\partial\Loss}{\partial\vect u}\right)^{\!\!\top}\;,}
  \qquad
  \text{equivalently } \mat A^\top \vect\eta
    = (\partial\Loss/\partial\vect u)^\top,
  \label{eq:adjoint}
\end{equation}
which costs a single linear solve. Then
$(\partial\Loss/\partial\vect u)\,\mat A^{-1}
= \vect\eta^\top$, and
\begin{equation}
  \boxed{\;
  \frac{\dd\Loss}{\dd\pts}
  \;=\;
  \vect\eta^\top \frac{\partial\vect b}{\partial\pts}
  - \vect\eta^\top \frac{\partial\mat A}{\partial\pts}\,\vect u
  + \frac{\partial\Loss}{\partial\pts}.
  \;}
  \label{eq:dLdp}
\end{equation}
This is the central identity. Every contribution to the gradient is
some piece of $\partial\mat A/\partial\pts$ or
$\partial\vect b/\partial\pts$ contracted against $\vect\eta$ and
$\vect u$, plus the explicit term.

\subsection{Decomposition of \texorpdfstring{$\partial\mat A/\partial\pts$}{dA/dp}}

The assembled $\mat A$ depends on $\pts$ through several intermediate
quantities. We unpack the chain rule, organizing the result by which
component of the assembly responds:
\begin{equation}
  \frac{\partial\mat A}{\partial\pts}
  \;=\;
  \sum_{\mathcal L \in \{\partial_{xx}, \partial_{yy},
       \partial_{xy}, \partial_x, \partial_y\}}
       \frac{\partial\mat A}{\partial \mat W^{\mathcal L}}
       \frac{\partial \mat W^{\mathcal L}}{\partial \pts}
  \;+\;
  \frac{\partial\mat A}{\partial\vect n}
  \frac{\partial\vect n}{\partial\pts}.
  \label{eq:dAdp-decomp}
\end{equation}
The first sum (over weight matrices) is the part that ``every PDE
has'': it captures how the operators change as the points move.
The second term (over normals) is specific to BCs that
re-orient under deformation; it is the Phase D Level 2
contribution and is zero whenever the normals are frozen.

\subsection{Decomposition of \texorpdfstring{$\partial\vect b/\partial\pts$}{db/dp}}

Symmetrically, $\vect b$ depends on $\pts$ through traction values
on Neumann rows and (in principle) through body forces on interior
rows:
\begin{equation}
  \frac{\partial\vect b}{\partial\pts}
  \;=\;
  \frac{\partial\vect b}{\partial\vect t}
  \frac{\partial\vect t}{\partial\pts}
  \;+\;
  \frac{\partial\vect b}{\partial\vect f}
  \frac{\partial\vect f}{\partial\pts}.
  \label{eq:dbdp-decomp}
\end{equation}
The first term is Phase D Level 1 (shape-dependent dead loads).
The second is identically zero in our cantilever example
(body-force-free) but the framework accommodates it.

\subsection{The five computational steps}

Combining~\eqref{eq:dLdp} with the decompositions~\eqref{eq:dAdp-decomp}
and~\eqref{eq:dbdp-decomp} yields the explicit form
\begin{equation}
  \begin{split}
  \frac{\dd\Loss}{\dd\pts}
  \;=\;\;&
  \underbrace{
    \vect\eta^\top \frac{\partial\vect b}{\partial\vect t}
    \frac{\partial\vect t}{\partial\pts}
    \;+\;
    \vect\eta^\top \frac{\partial\vect b}{\partial\vect f}
    \frac{\partial\vect f}{\partial\pts}
  }_{\text{Step 5b: }\partial\vect b/\partial\pts}
  \\
  &\;\;-\;\;
  \underbrace{
    \sum_{\mathcal L}
    \vect\eta^\top \frac{\partial\mat A}{\partial\mat W^{\mathcal L}}
    \frac{\partial\mat W^{\mathcal L}}{\partial\pts}\,\vect u
  }_{\text{Steps 3+4: } \partial\mat W/\partial\pts}
  \;\;-\;\;
  \underbrace{
    \vect\eta^\top \frac{\partial\mat A}{\partial\vect n}
    \frac{\partial\vect n}{\partial\pts}\,\vect u
  }_{\text{Step 5c: } \partial\vect n/\partial\pts}
  \;\;+\;\;
  \underbrace{
    \frac{\partial\Loss}{\partial\pts}.
  }_{\text{Step 5a: explicit}}
  \end{split}
  \label{eq:five-pieces}
\end{equation}
This decomposes into the five computational steps that the
implementation organizes:
\begin{enumerate}
  \item \textbf{Forward solve}: $\vect u = \mat A^{-1}\vect b$,
        caching weight matrices and their RBF stencil factorizations.
  \item \textbf{Adjoint solve}: $\mat A^\top \vect\eta
        = (\partial\Loss/\partial\vect u)^\top$, single linear solve.
  \item \textbf{Extract $\Delta\mat W$}: convert
        $-\vect\eta^\top (\partial\mat A/\partial\mat W^{\mathcal L})
        \vect u$ into a per-entry $\Delta\mat W^{\mathcal L}$ for each
        operator, treating sparsity and per-row coefficients
        analytically.
  \item \textbf{Backward through RBF stencils}: convert each
        $\Delta\mat W^{\mathcal L}$ into a contribution to
        $\Delta\pts$ by calling
        \texttt{\_pullback\_weights!}.
  \item \textbf{External contributions}: add the explicit
        $\partial\Loss/\partial\pts$ (Step 5a), the
        $\partial\vect b/\partial\pts$ piece (Step 5b), and the
        Phase D Level 2 $\partial\vect n/\partial\pts$ piece
        (Step 5c).
\end{enumerate}
We now detail each step.

\section{The five-step pipeline in detail}
\label{sec:five-steps}

\subsection{Step 1 --- forward solve and caching}

Given $\pts$, build the five RBF weight matrices needed for elasticity
($\mat W^{\partial_{xx}}, \mat W^{\partial_{yy}},
\mat W^{\partial_{xy}}$ on all points; $\mat W^{\partial_x},
\mat W^{\partial_y}$ on the Neumann sub-cloud). For each weight
matrix, also retain the per-row \emph{cache} returned by
\texttt{\_build\_weights\_and\_cache} ---
this is the saddle-point factorization used to produce the row, and
it is what \texttt{\_pullback\_weights!} reuses for the backward pass.

Assemble $\mat A$ by combining the interior block~\eqref{eq:A-int}
with Dirichlet identity rows~\eqref{eq:dirichlet} and Neumann
rows~\eqref{eq:neumann-row}. Assemble $\vect b$ from prescribed
Dirichlet values, traction values $\vect t$, and (if present) body
forces.

Factor $\mat A$ once (LU) and solve $\vect u = \mat A^{-1}\vect b$.
The same factorization will be reused for Step 2.

\subsection{Step 2 --- adjoint solve}

Form $(\partial\Loss/\partial\vect u)^\top$ analytically (cheap for
the typical losses --- $2\vect u$ for $\Loss = \norm{\vect u}^2$,
$\vect b$ for compliance $\Loss = \vect b^\top \vect u$, etc.). Solve
$\mat A^\top \vect\eta = (\partial\Loss/\partial\vect u)^\top$ using
the cached factorization of $\mat A$. Cost: one back-substitution.

For compliance specifically, $\vect\eta = \mat A^{-\top} \vect b
= \vect u$ if $\mat A$ is symmetric --- the so-called
\emph{self-adjoint} case. Our $\mat A$ is symmetric for the standard
elasticity row family (it can be made so by symmetrizing the
Neumann rows), but we run the explicit adjoint solve anyway to
avoid relying on this property when the formulation is extended.

\subsection{Step 3 --- extract \texorpdfstring{$\Delta\mat W$ from $\Delta\mat A$}{Delta W from Delta A}}
\label{sec:step3}

This is the most code-dense step. The goal is to translate
\begin{equation}
  \Delta\mat A
  \;=\; -\vect\eta\,\vect u^\top
  \quad \in \reals^{2N \times 2N}
\end{equation}
into per-operator sensitivities $\Delta\mat W^{\mathcal L}
\in \reals^{N \times N}$ that capture
$-\vect\eta^\top(\partial\mat A/\partial\mat W^{\mathcal L})\vect u$
in a single expression per nonzero entry. The decomposition is
done analytically using the elasticity coefficient structure
in~\eqref{eq:A-int} (interior) and~\eqref{eq:neumann-row} (Neumann).

\subsubsection{Interior rows}

Each interior row family is governed by~\eqref{eq:A-int}. For a pair
of points $(i, j)$ with $j$ in the stencil of $i$, the four
$2\times 2$ block entries are
\begin{align}
  \mat A[i,j]      &= c_1 \mat W^{\partial_{xx}}_{ij} + c_2 \mat W^{\partial_{yy}}_{ij},
  & \mat A[i, j{+}N]   &= c_3 \mat W^{\partial_{xy}}_{ij}, \\
  \mat A[i{+}N, j] &= c_3 \mat W^{\partial_{xy}}_{ij},
  & \mat A[i{+}N, j{+}N] &= c_2 \mat W^{\partial_{xx}}_{ij} + c_1 \mat W^{\partial_{yy}}_{ij}.
\end{align}
Setting $t_{11} = -\eta_i u_j$, $t_{22} = -\eta_{i+N} u_{j+N}$,
$t_{12} = -\eta_i u_{j+N}$, $t_{21} = -\eta_{i+N} u_j$ (the four
$2{\times}2$ entries of $\Delta\mat A[i,j;\, i+N, j+N]$), the
back-projection onto the operators is
\begin{equation}
  \boxed{\;
  \begin{aligned}
    \Delta\mat W^{\partial_{xx}}_{ij} \;&=\; c_1 t_{11} + c_2 t_{22}, \\
    \Delta\mat W^{\partial_{yy}}_{ij} \;&=\; c_2 t_{11} + c_1 t_{22}, \\
    \Delta\mat W^{\partial_{xy}}_{ij} \;&=\; c_3 (t_{12} + t_{21}).
  \end{aligned}
  \;}
  \label{eq:interior-extract}
\end{equation}
This is the \texttt{extract\_weight\_sensitivities\_elasticity!}
routine in \texttt{src/optimization/manual\_adjoint.jl}. It walks
the shared sparsity pattern of the three weight matrices once and
populates all three $\Delta\mat W$ in a single pass.

\paragraph{Per-DOF active mask.}
At rows that are overwritten with Dirichlet identity (where
$\partial\mat A/\partial\mat W = 0$ for that row), the
contribution must be gated to zero so we do not accidentally pick
up the identity-row sensitivities. We carry an
\texttt{active::Vector\{Bool\}} of length $2N$ and apply the
gating per-DOF:
\begin{equation}
  t_{ab} \;=\; \begin{cases}
    \text{(formula above)} & \text{if active}[i\,(+N\,\text{if}\,a{=}2)] \\
    0 & \text{otherwise}.
  \end{cases}
\end{equation}
The per-DOF granularity (not per-point) matters: a corner can be
Dirichlet in $x$ but Neumann in $y$, and the mask correctly
suppresses only the Dirichlet half.

\paragraph{Skipping Neumann rows in the interior loop.}
The interior block~\eqref{eq:A-int} is not the equation actually
imposed on Neumann rows after overwriting in~\eqref{eq:neumann-row}.
The Step 3 loop therefore restricts the interior extraction to
rows where the interior assembly is in effect (i.e.\ rows not in
$\Gamma_D$ and not in $\Gamma_N$), via an
\texttt{interior\_rows::BitVector} parameter.

\subsubsection{Neumann rows}
\label{sec:step3-neumann}

For each Neumann point with local index $i_{\text{loc}}$ and global
index $g$, the Neumann row family~\eqref{eq:neumann-row} writes each
$\mat A[g,\,\cdot]$ and $\mat A[g{+}N,\,\cdot]$ entry as a linear
combination of $\mat W^{\partial_x}[i_{\text{loc}},\,\cdot]$ and
$\mat W^{\partial_y}[i_{\text{loc}},\,\cdot]$ with coefficients
that depend on $(n_x, n_y, \lambda^\star, \mu)$. Storing these
coefficients in a flat array \texttt{coeffs}, the
\texttt{TractionLayout} struct catalogues, for each entry, its
$(c_x, c_y)$ pair and the $(w_{\text{col}})$ index in the weight
matrices. Given the Step 3 base assignment $\Delta a = -\eta_g u_c$,
the back-projection onto the weight matrices is simply
\begin{equation}
  \boxed{\;
  \Delta\mat W^{\partial_x}_{i_{\text{loc}},\, w_{\text{col}}}
    \mathrel{+}= c_x \Delta a, \qquad
  \Delta\mat W^{\partial_y}_{i_{\text{loc}},\, w_{\text{col}}}
    \mathrel{+}= c_y \Delta a.
  \;}
\end{equation}
This is \texttt{extract\_neumann\_sensitivities!}. It is gated by
the same \texttt{active} mask and accumulates additively into the
$\Delta\mat W$ already populated by the interior extraction.

\subsection{Step 4 --- backward through the RBF stencils}
\label{sec:step4}

Given $\Delta\mat W^{\mathcal L}$ for each operator, we need
$\Delta\pts$ contributions such that
$\partial\mat W^{\mathcal L}/\partial\pts \cdot \Delta\pts^\top
= \Delta\mat W^{\mathcal L}$ in the standard pullback sense. This
is provided by \texttt{\_pullback\_weights!} in
\texttt{RadialBasisFunctions.jl}, which we call once per operator.
The routine reuses the saddle-point factorization cached in Step 1
and runs purely as Julia linear algebra --- no automatic
differentiation framework is involved.

For each operator we call
\begin{equation}
  \Delta\pts \;\mathrel{+}=\;
  \texttt{\_pullback\_weights!}\big(
    \Delta\mat W^{\mathcal L},\;
    \mat W^{\mathcal L},\;
    \text{cache},\;
    \pts,\;
    \text{eval\_pts},\;
    \text{adjl},\;
    \text{basis},\;
    \mathcal L
  \big),
\end{equation}
where \texttt{eval\_pts} is the full point cloud for the interior
operators and the Neumann sub-cloud for the Neumann operators
(with an \texttt{eval\_offset} array to map back to global indices
for the Neumann case). The cost per operator is
$O(N k^3)$ for the per-row cache replay, dominated by the cube of
stencil size.

\subsection{Step 5 --- external contributions}

The three pieces collected in Step 5 are independent of the
$\Delta\mat W$ pipeline and are added directly to $\Delta\pts$.

\subsubsection{Step 5a --- the explicit
\texorpdfstring{$\partial\Loss/\partial\pts$}{dL/dp}}
\label{sec:step5-explicit}

The loss may depend on $\pts$ directly through geometric quantities.
The canonical case is the area penalty
$\Loss_{\text{geom}}(\pts) = \rho_{\text{pen}}(V(\pts) - V_0)^2$.
With the shoelace formula on the boundary loop,
$\partial V/\partial \pts_i$ is sparse (nonzero only on boundary
points) and analytic. This term is added in the example via the
helper \texttt{polygon\_area\_grad!}.

\subsubsection{Step 5b ---
\texorpdfstring{$\partial\vect b/\partial\pts$ (Phase D Level 1)}
{db/dp (Phase D L1)}}
\label{sec:step5-load}

When the traction $\vect t$ depends on $\pts$ --- the
\emph{dead-load} case --- the contribution
$\vect\eta^\top\,\partial\vect b/\partial\pts$ is non-zero. The
dependence can take two qualitatively different forms, and the
machinery is set up to handle both:
\begin{description}[leftmargin=2em,style=nextline]
  \item[\textbf{Direct local dependence.}]
    The traction at Neumann point $g$ is a function only of $\pts_g$
    itself: $\vect t_g = \vect t(\pts_g)$. A typical case is a
    profile prescribed by location, e.g.\ a parabolic shear
    $t_y(y) = P(D^2 - y^2)/(2I)$ depending on the local $y$-coord.
    The Jacobian
    $\mat J^{(i_{\text{loc}})} \equiv \partial\vect t_g/\partial\pts_g$
    is a single $2\times 2$ matrix per Neumann point.
  \item[\textbf{Indirect non-local dependence.}]
    The traction at Neumann point $g$ depends not only on $\pts_g$
    but on the coordinates of \emph{other} boundary points through
    an intermediate geometric quantity. For example, in the
    cantilever the same parabolic-shear formula with a
    physically-meaningful effective half-depth
    $D_{\text{eff}} = (y_{\text{top corner}} - y_{\text{bot corner}})/2$
    induces dependence of every right-edge Neumann traction on the
    two corner $y$-coordinates via
    $\partial \vect t/\partial D_{\text{eff}}
     \cdot \partial D_{\text{eff}}/\partial \pts_{\text{corner}}$.
    The Jacobian
    $\mat J^{(i_{\text{loc}}, q)} \equiv \partial\vect t_g/\partial\pts_q$
    is now in general non-zero for several $q \ne g$.
\end{description}
In both cases the gradient contribution is the same chain-rule
product; the only difference is how many destinations the gradient
flows to:
\begin{equation}
  \boxed{\;
  \Delta\pts_q \;\mathrel{+}=\;
  \begin{pmatrix} \eta_g \\ \eta_{g+N} \end{pmatrix}^{\!\top}\!
  \mat J^{(i_{\text{loc}}, q)},
  \qquad
  \mat J^{(i_{\text{loc}}, q)}_{a,b}
    = \frac{\partial t_a}{\partial p_{q,b}}.\;}
  \label{eq:l1-eta-side}
\end{equation}
The direct case is handled by
\texttt{extract\_load\_sensitivities!} called with a per-Neumann-point
$2\times 2$ Jacobian. The indirect case is handled by a thin
companion routine that, for each Neumann point whose traction
depends on a non-local quantity $Q(\pts)$, contracts
$(\partial t/\partial Q)$ with $\eta_g, \eta_{g+N}$ and distributes
the result through the (typically sparse) Jacobian
$\partial Q/\partial\pts$. The $D_{\text{eff}}$ case above is one
example; a follower pressure
$\vect t = -p\,\vect n$ would be another, with the gradient
flowing through both the polyline-normal Jacobian
of Section~\ref{sec:l2} and a $\partial p/\partial\pts$ if the
pressure itself is shape-coupled.

\paragraph{Why this is not the only $b$-side contribution.}
The compliance loss $\Loss = \vect b^\top \vect u$ depends on
$\vect b$ \emph{explicitly}: $\partial\Loss/\partial\vect b
= \vect u^\top$. The total $b$-side sensitivity is therefore
\begin{equation}
  \frac{\dd\Loss}{\dd\pts}\Big|_b
  \;=\;
  \vect\eta^\top \frac{\partial\vect b}{\partial\pts}
  \;+\;
  \frac{\partial\Loss}{\partial\vect b}\,
  \frac{\partial\vect b}{\partial\pts}
  \;=\;
  (\vect\eta + \vect u)^\top \frac{\partial\vect b}{\partial\pts}
\end{equation}
for compliance. The first term is the $\vect\eta$-side contribution
and is computed inside \texttt{shape\_gradient} when the traction
Jacobian is supplied. The second is the $\vect u$-side
contribution and must be added by the caller, by a \emph{second}
call to the same load-sensitivity routine with $\vect u$ in place
of $\vect\eta$. Missing the second term introduces a $\sim 40\%$
FD error on a compliance loss --- a hazard explicitly flagged in
\texttt{CLAUDE.md}.

\paragraph{A note on the adjoint output.}
For losses with non-trivial $\partial\Loss/\partial\vect b$, the
caller needs $\vect\eta$ \emph{and} $\vect u$ available outside
\texttt{shape\_gradient} so it can run the second
load-sensitivity contraction. The pipeline therefore returns
$\vect\eta$ alongside $\vect u$ as part of the
\texttt{shape\_gradient} result, avoiding a redundant adjoint
solve at the caller.

\subsubsection{Step 5c --- \texorpdfstring{$\partial\vect n/\partial\pts$}{dn/dp}
(Phase D Level 2)}
\label{sec:step5-normals}

This is the term that closes the differentiability of the
Neumann assembly when the normals are shape-dependent. We give
it a dedicated section (Section~\ref{sec:l2}) because the
derivation is substantial.

\section{Differentiable Neumann data}
\label{sec:boundary-diff}

\subsection{Phase D Level 1 --- shape-dependent traction values}

The case where $\vect t$ depends on $\pts$ was covered in
Section~\ref{sec:step5-load}. The implementation surface is small ---
a kwarg \texttt{traction\_jacobians} on \texttt{shape\_gradient},
plus the helper \texttt{extract\_load\_sensitivities!} ---
because the analytic content (the per-point $2\times 2$
$\partial\vect t/\partial\pts_g$ matrix) is supplied by the
caller, and the framework just does the contraction with
$\vect\eta$ (and, if needed, with $\vect u$). The framework is
backward-compatible: omitting \texttt{traction\_jacobians} reduces
to the frozen-load case used by earlier phases.

\subsection{Phase D Level 2 --- differentiable normals}
\label{sec:l2}

\subsubsection{Motivation}

The normals $\vect n$ enter the Neumann row coefficients
in~\eqref{eq:neumann-row}. If $\vect n$ is frozen at a reference
value, $\partial\mat A/\partial\vect n \cdot \partial\vect n/\partial\pts
= 0$ and the Neumann assembly contributes to the gradient only
through $\partial\mat W^{\partial_x}/\partial\pts$ and
$\partial\mat W^{\partial_y}/\partial\pts$. This is sufficient
when boundary points are constrained to move along their reference
edge (so their normals do not rotate), but fails for any geometry
where boundary points can move freely --- in particular, for
corner DOFs and for free-form boundaries.

\subsubsection{Chord-formula vertex normal}

Let $\mathcal P = (P_1, P_2, \ldots, P_M)$ be the CCW-ordered closed
polyline of boundary vertices, with cyclic indexing. For Neumann
point $i_{\text{loc}}$ at position $k$ in $\mathcal P$ we define
\begin{equation}
  \vect n_{i_{\text{loc}}}(\pts)
  \;=\;
  \mat R(-\pi/2)\,
  \frac{\pts_{P_{k+1}} - \pts_{P_{k-1}}}{\norm{\pts_{P_{k+1}} - \pts_{P_{k-1}}}},
  \qquad
  \mat R(-\pi/2) = \begin{pmatrix} 0 & 1 \\ -1 & 0 \end{pmatrix}.
  \label{eq:chord-normal}
\end{equation}
At an edge midpoint where $P_{k\pm 1}$ are collinear with $P_k$
this reduces to the canonical edge normal; at a polygon corner it
yields a length-weighted bisector, biased toward the longer of the
two adjacent edges. We discuss the corner subtlety in
Section~\ref{sec:corner-fix}.

\subsubsection{Closed-form sparse Jacobian}

With $\vect d \equiv \pts_{P_{k+1}}-\pts_{P_{k-1}}$,
$\ell = \norm{\vect d}$, and unit tangent
$\vect t = \vect d/\ell = (t_x, t_y)$, direct differentiation
of~\eqref{eq:chord-normal} gives
\begin{equation}
  \frac{\partial\vect n}{\partial\vect d}
  \;=\;
  \frac{1}{\ell}
  \begin{pmatrix}
    -t_x t_y & \phantom{-}t_x^2 \\
    -t_y^2  & \phantom{-}t_x t_y
  \end{pmatrix},
  \qquad
  \frac{\partial\vect d}{\partial \pts_{P_{k+1}}} = +\mat I, \quad
  \frac{\partial\vect d}{\partial \pts_{P_{k-1}}} = -\mat I.
  \label{eq:chord-jac}
\end{equation}
The Jacobian of $\vect n_{i_{\text{loc}}}$ with respect to all
boundary coordinates is sparse with nonzeros at exactly two
vertices, $P_{k-1}$ and $P_{k+1}$. We package this in a
\texttt{NormalJacobian} struct storing the four $\reals^2$ blocks
$\partial n_a/\partial \pts_{P_{k\pm 1}}$ for $a \in \{x, y\}$.
A standalone unit test
(\texttt{examples/test\_polyline\_normals.jl}) verifies
\eqref{eq:chord-jac} against \texttt{FiniteDifferences.central\_fdm(5,1)}
with worst-case relative error $3.9\times 10^{-12}$ on a perturbed
rectangle.

\subsubsection{Adjoint contribution from \texorpdfstring{$\partial\vect n/\partial\pts$}{dn/dp}}

Returning to~\eqref{eq:five-pieces}, the L2 contribution is
$- \vect\eta^\top (\partial\mat A/\partial\vect n)
   (\partial\vect n/\partial\pts) \vect u$. Each Neumann row $g$
contributes
\begin{equation}
  -\eta_g \sum_c
   \frac{\partial\mat A[g,c]}{\partial\vect n}
   \frac{\partial\vect n_{i_{\text{loc}}}}{\partial\pts_q}\,
   u_c.
\end{equation}
Because each Neumann row entry~\eqref{eq:neumann-row} is a linear
function of $(n_x, n_y)$, the partial derivatives
$\partial\mat A[g,c]/\partial n_a$ are constants drawn from
$\{0, \lambda^\star, \mu, \lambda^\star{+}2\mu\}$ multiplied by
either $\mat W^{\partial_x}_{i_{\text{loc}},c'}$ or
$\mat W^{\partial_y}_{i_{\text{loc}},c'}$. Pre-contracting over
$c$ yields scalar coefficients $r^{(g)}_a$:
\begin{equation}
  r^{(g)}_a
  \;\equiv\;
  \sum_{p \in \text{entries}(g)}
    \frac{\partial \alpha_p}{\partial n_a}\,
    \mat W^{\mathcal L_p}[i_{\text{loc}}, w_p]\, u_{c_p},
  \qquad a \in \{x, y\}.
  \label{eq:row-contraction}
\end{equation}
The L2 contribution to $\Delta\pts_q$ is then
\begin{equation}
  \boxed{\;
  \Delta\pts_q \;\mathrel{+}=\;
  -\!\!\!\sum_{a \in \{x,y\}}
   \big(\eta_g\, r^{(g)}_a + \eta_{g+N}\, r^{(g+N)}_a\big)\,
   \frac{\partial n_a}{\partial \pts_q},
   \;}
  \label{eq:l2-distribute}
\end{equation}
nonzero only for $q \in \{P_{k-1}, P_{k+1}\}$ via the sparse
\texttt{NormalJacobian}. This is the routine
\texttt{extract\_normal\_sensitivities!}. It is invoked from
\texttt{shape\_gradient} when the caller supplies the
\texttt{normal\_jacobians} kwarg; default \texttt{nothing}
preserves the original frozen-normal behavior.

\subsubsection{The corner pathology}
\label{sec:corner-fix}

For a 90$^\circ$ corner with adjacent edge segments of lengths
$\Delta x \ne \Delta y$, the chord normal~\eqref{eq:chord-normal}
is \emph{not} the angular bisector of the two adjacent edge
normals. At the cantilever's top-right corner with
$\Delta x = 1.0, \Delta y = 0.5$, the chord formula gives
$\vect n = (1, 2)/\sqrt 5 \approx (0.447, 0.894)$, biased toward
the top-edge normal. Imposing $\vect\sigma\!\cdot\!\vect n = \vect 0$
at the corner with this normal gives the linear combination
$\sigma_{xx} + 2\sigma_{xy} = 0$ and $\sigma_{xy} + 2\sigma_{yy} = 0$,
which is misaligned with the beam-natural condition
$\sigma_{xx} = 0$ at the free end.

In the cantilever example we resolve this pragmatically by
overriding the corner normal to the canonical $(1,0)$ and zeroing
the corresponding \texttt{NormalJacobian} entries. The corner
\emph{coordinates} remain free design variables; the L2
sensitivity still flows through the corner because the corner
appears as one of the $(P_{k-1}, P_{k+1})$ neighbors of its
adjacent edge midpoints, and those midpoints' chord normals rotate
when the corner moves. The general handling of corners ---
selecting between the chord formula, the angular bisector, or a
physics-aware convention --- belongs in the geometry kernel and is
discussed in Section~\ref{sec:openwork}.

\section{Sensitivity filtering on the boundary loop (superseded)}
\label{sec:filter}

\emph{This section records the regularization first used --- in the cantilever
example (Section~\ref{sec:cantilever}) --- and later superseded.}
Section~\ref{sec:plate} shows that filtering the gradient is the \emph{wrong} cure
for RBF-FD boundary noise: the noise is broadband and overlaps the signal, so no
boundary filter separates them faithfully; the structural cure is the design-space
contraction~\eqref{eq:design-contraction}. The filter definitions are kept here
because the cantilever results reference them.

\subsection{Why we filtered}

The raw discrete gradient $\Delta\pts^{\text{raw}} \equiv \dd\Loss/\dd\pts$ is the
exact gradient of the discrete loss, but on a free-node design it is contaminated by
a near-Nyquist mode along the boundary (FD-validated to $10^{-7}$, yet rough):
unfiltered $\ell^2$ descent amplifies it into a checkerboard.

\subsection{The 3-point Nyquist filter}

The simplest mitigation is a single-pass 3-point average along the
CCW boundary loop:
\begin{equation}
  g_k^{\text{smooth}}
  \;=\; \alpha\, g_{k-1} + (1 - 2\alpha)\, g_k + \alpha\, g_{k+1},
  \qquad \alpha = \tfrac{1}{4}.
\end{equation}
The transfer function for a periodic sinusoid of wavenumber
$\omega$ per segment is
\begin{equation}
  H_3(\omega) \;=\; 1 - 4\alpha\sin^2(\omega/2)
  \;=\; \cos^2(\omega/2)\quad\text{at }\alpha=\tfrac14.
\end{equation}
This is a strict zero at $\omega = \pi$ (Nyquist mode killed) but
leaves intermediate frequencies largely intact:
$H_3(\pi/2) = \tfrac12$, $H_3(\pi/4) \approx 0.85$. The optimizer
can therefore still pump energy into modes around half-Nyquist and
below.

\subsection{The Helmholtz-PDE filter}

To suppress medium-wavelength modes we solve, instead, an elliptic
PDE on the closed boundary loop:
\begin{equation}
  (\mat I - r^2\,\nabla^2_{\text{loop}})\, \vect g^{\text{smooth}}
  \;=\; \vect g^{\text{raw}},
  \label{eq:helmholtz}
\end{equation}
where $\nabla^2_{\text{loop}}$ is the standard 1D periodic Laplacian
on the loop, stencil $[-1, 2, -1]$ at each vertex with its cyclic
neighbors. The parameter $r$ has units of boundary segments and
sets the length scale of attenuation. Frozen boundary entries are
imposed as identity rows in the system so the gradient is
attenuated as it approaches a clamped (Dirichlet) region.

The Fourier-mode multiplier is
\begin{equation}
  H_H(\omega; r) \;=\; \frac{1}{1 + 4r^2\sin^2(\omega/2)}.
  \label{eq:hh-transfer}
\end{equation}
Selected values:

\begin{table}[h]
\centering
\caption{Filter transfer at representative wavenumbers.
$H_3$ is the 3-point average; $H_H(\cdot;r)$ is the Helmholtz
filter~\eqref{eq:hh-transfer}.}
\label{tab:filter-transfer}
\begin{tabular}{lccc}
\toprule
$\omega$ & $H_3$ & $H_H(r=1)$ & $H_H(r=2)$ \\
\midrule
$0$           (DC)            & 1.00 & 1.00 & 1.00 \\
$\pi/4$       (period 8)      & 0.85 & 0.62 & 0.29 \\
$\pi/2$       (period 4)      & 0.50 & 0.33 & 0.09 \\
$\pi$         (Nyquist)       & 0.00 & 0.20 & 0.06 \\
\bottomrule
\end{tabular}
\end{table}

\subsection{Implementation}

The discrete system is a cyclic tridiagonal of size $M$ (boundary-loop vertices),
solved in $O(M)$ (Thomas $+$ Sherman--Morrison; a dense solve is microseconds at
$M\sim24$). See \texttt{helmholtz\_filter\_along\_boundary!} in the Phase-3
cantilever example.

\section{Verification methodology}
\label{sec:verification}

The discrete-adjoint pipeline is checked at every phase by
comparing the analytic gradient to a fifth-order central
finite-difference reference. The validation is built into each of
the standalone Phase A and Phase B examples and into the
illustrative Phase 3 cantilever:
\begin{equation}
  \text{rel\_err}(\pts)
  \;\equiv\;
  \frac{\norm{\nabla_{\pts}\Loss\,\big|_{\text{AD}}
              - \nabla_{\pts}\Loss\,\big|_{\text{FD}}}_2}{\norm{\nabla_{\pts}\Loss\,\big|_{\text{FD}}}_2},
\end{equation}
with the FD reference computed by
\texttt{FiniteDifferences.central\_fdm(5,1)}. The targets are
machine precision: rel\_err $\sim 10^{-11}$ for Dirichlet-only
problems (Phase A), $\sim 10^{-7}$ for problems with traction BCs
(Phase B and beyond --- the lower precision is due to FD truncation
error on the larger relevant scale, not to the AD pipeline).

For the cantilever in Section~\ref{sec:cantilever} we additionally
probe two qualitative properties:
\begin{itemize}[leftmargin=2em]
  \item \textbf{Top/bottom symmetry.}\;
        The cantilever geometry and loading are symmetric under
        $y \to -y$. We match top-edge nodes to their bottom-edge
        mirrors by $x$-coordinate and compute
        $\max_i \lvert g_y(\text{top}_i) + g_y(\text{bot}_i)\rvert$.
        On the reference geometry this is $\sim 10^{-8}$, i.e.\
        symmetric to floating-point rounding.
  \item \textbf{Spectral content of the raw gradient.}\;
        We project the iter-0 gradient onto the alternating-sign
        mode $\phi_\pi$ of the boundary loop and check its weight.
        Globally this is essentially zero ($\sim 10^{-9}$); on
        each individual edge, however, the raw gradient is
        strongly Nyquist-aligned, by an amount the 3-point filter
        is designed to suppress.
\end{itemize}
Together these three diagnostics --- finite differences, symmetry,
and spectral content --- are sufficient to distinguish a true
implementation error (which would fail finite-differences) from a
regularization issue (which would fail spectral content but pass
finite differences) from a discretization symmetry break (which
would fail symmetry alone).

\section{Illustrative example: the cantilever}
\label{sec:cantilever}

\subsection{Setup}

We illustrate the pipeline end-to-end on a 2D cantilever. Domain
$[0, L] \times [-D, D]$ with $L = 8.0$, $D = 1.0$, modulus
$E = 10^7$, Poisson ratio $\nu = 0.3$, RBF basis $\phi(r) = r^3$
with augmenting polynomial degree 3 and stencil size $k = 35$.
We discretize with a 9$\times$5 regular grid ($N = 45$). The
loading is a parabolic shear traction on the right edge,
$t_y(y) = P(D^2 - y^2)/(2I)$ with $P = 10^3$, $I = 2D^3/3$, zero
on top and bottom. The left edge is clamped.

The loss is compliance plus a two-sided area penalty:
\begin{equation}
  \Loss(\pts) = \vect b(\pts)^\top \vect u(\pts)
    \;+\; \rho_{\text{pen}}\big(V(\pts) - V_0\big)^2,
  \qquad \rho_{\text{pen}} = 10^3.
\end{equation}
Design variables are the $y$-coordinates of the Neumann-boundary
points; the optimizer is plain gradient descent with Armijo
backtracking; the iteration budget is 40.

\subsection{End-to-end workflow}

At each iteration:
\begin{enumerate}
  \item Refresh the live traction values and (when L2 is active)
        the live polyline normals on \texttt{traction\_layout}.
  \item Compute the analytic gradient via \texttt{shape\_gradient}
        (Steps 1--4 + the Phase D L1 $\vect\eta$-side b-contribution).
  \item Add the $\vect u$-side Phase D L1 contribution
        (\texttt{extract\_load\_sensitivities!} with $\vect u$).
  \item Add the area-penalty gradient (Step 5a).
  \item Apply the sensitivity filter on the boundary loop, then
        mask to the free DOFs.
  \item Armijo backtracking line search; update $\pts$.
\end{enumerate}

\subsection{Snapshot results}

Numerical results are reported here as a snapshot of the
pipeline's qualitative behavior during active development. The
\emph{absolute} compliance figures may shift as the cantilever
example is refined --- in particular, an artefact in the
tip-shrinkage gradient sign was masked by an earlier corner-freeze
convention and is being repaired at the time of writing. The
\emph{qualitative} contrasts between configurations (effect of L2
on the corner DOF, effect of filter choice on the boundary
smoothness) are robust and survive that refinement.

\begin{table}[h]
\centering
\caption{Snapshot of three configurations on the cantilever
benchmark, 40 iterations each. All configurations pass the
finite-difference validation at iter 0 to
$\text{rel\_err} \sim 10^{-7}$.}
\label{tab:results}
\begin{tabular}{lrr}
\toprule
Configuration & $\Delta C / C_0$ & boundary appearance \\
\midrule
(A) corners frozen, 3-point filter           & $-47.6\%$ & mild waviness \\
(B) corners free, 3-point filter             & $-52.1\%$ & visible long-wavelength undulation \\
(C) corners free, Helmholtz $r{=}1$ filter   & $-46.6\%$ & visibly smooth \\
\bottomrule
\end{tabular}
\end{table}

\begin{figure}[h]
\centering
\includegraphics[width=0.95\linewidth]{phase3_corners_frozen.png}
\caption{Configuration (A): corners frozen, 3-point filter.}
\label{fig:A}
\end{figure}

\begin{figure}[h]
\centering
\includegraphics[width=0.95\linewidth]{phase3_l2_3point.png}
\caption{Configuration (B): corners free (Phase D L2 active),
3-point filter. The optimizer extracts additional compliance
reduction by exploiting medium-wavelength modes the 3-point filter
leaves unattenuated, producing visible boundary undulations.}
\label{fig:B}
\end{figure}

\begin{figure}[h]
\centering
\includegraphics[width=0.95\linewidth]{phase3_l2_helmholtz_r1.png}
\caption{Configuration (C): corners free, Helmholtz $r{=}1$ filter.
The boundary is now visibly smooth. The compliance reduction is
slightly lower than (B) because the optimizer is restricted to the
smooth design subspace; this is the honest optimum under physical
regularization rather than the apparent optimum that exploits
discretization noise.}
\label{fig:C}
\end{figure}

\subsection{What the three configurations demonstrate}

Despite the noted ongoing refinement of the cantilever example, the
contrast between (A)--(C) carries information that is robust:

\begin{enumerate}[label=(\arabic*)]
  \item Going from (A) to (B) demonstrates that the Phase D L2
        machinery is doing real algorithmic work: unfreezing the
        corner DOFs gives the optimizer more design freedom, and
        the gradient through the chord-formula normals is correct
        (FD-validated).
  \item Going from (B) to (C) demonstrates that the 3-point
        filter is insufficient for shape regularization in the
        presence of richer design freedom: the optimizer exploits
        unfiltered modes, and the Helmholtz filter is necessary to
        recover a smooth-design optimum.
  \item Configuration (C) is the appropriate baseline against
        which to compare future algorithmic improvements (adaptive
        $r$, augmented-Lagrangian area constraint, modern
        optimizer): the regularization is honest and the
        compliance number is not inflated by noise exploitation.
\end{enumerate}

\section{The plate with a hole: from gradient filtering to a
design-space and re-meshing architecture}
\label{sec:plate}

The cantilever validated the adjoint but left a question the filter of
Section~\ref{sec:filter} only papered over: \emph{why} does the raw per-node
gradient need filtering at all, and is filtering the right cure? The
plate-with-hole problem --- a square plate in biaxial tension with an elliptical
hole, optimized toward the circular hole that minimizes compliance --- answered
both, and in doing so reorganized the optimizer around a design parameterization
and an actively re-meshed cloud. This section records that turn; the granular
record is the companion log \texttt{boundary\_gradient\_noise.md}.

\begin{figure}[h]
\centering
\includegraphics[width=0.6\linewidth]{fig_setup.png}
\caption{The plate-with-hole discretization: the interior cloud, the outer Neumann
frame under biaxial traction $\sigma_\infty$, the free (traction-free) hole
boundary, and the three Dirichlet pins that remove rigid-body modes.}
\label{fig:plate-setup}
\end{figure}

\subsection{The symptom and the diagnosis:
\texorpdfstring{$\partial\mat W/\partial\pts$}{dW/dp}}
\label{sec:plate-noise}

On grids fine enough to resolve the hole, the per-node $\ell^2$ shape gradient is
dominated by a near-Nyquist sawtooth: the gradient spectrum \emph{rises} toward
the highest representable boundary mode, so a near-Nyquist wiggle carries
$\sim\!40\times$ the gradient of the circle-seeking $\cos 2\theta$ mode. The
gradient is not wrong --- it is finite-difference validated mode by mode to
$10^{-6}$, the near-Nyquist modes included --- it is genuinely \emph{rough}.

Fixed-cloud probes localized the mechanism. With the PDE removed and replaced by
smooth analytic fields $f,\eta$, the surrogate $S = \eta^\top\mat W f$
differentiated with respect to a boundary node's radial position still produces a
broadband, Nyquist-rich $\partial S/\partial r_j$. The roughness is therefore the
\emph{geometry sensitivity of the RBF-FD differentiation weights themselves},
$\partial\mat W/\partial\pts$, independent of the physics (it reproduces on a
thermal Laplace problem). It grows with the operator's derivative order (the
$q{=}2$ elasticity/Laplacian operators are $\sim\!46\times$ worse than the $q{=}1$
Neumann rows, because the pull-back of a second-order operator probes the third
derivative of the $r^3$ kernel), worsens with kernel order (a conditioning
effect), and sharpens under refinement. Crucially it is \emph{broadband} and
overlaps the physical signal in frequency.

\begin{figure}[h]
\centering
\includegraphics[width=0.95\linewidth]{fig_noise.png}
\caption{The boundary gradient noise. \emph{Left:} the per-node radial shape
gradient $\partial C/\partial r_j$ walked around the hole --- a near-Nyquist
sawtooth. \emph{Right:} its Fourier amplitude spectrum \emph{rises} toward the
highest representable mode; the circle-seeking $m{=}2$ mode (green) sits far below
the high-mode band. The gradient is FD-exact yet genuinely rough.}
\label{fig:plate-noise}
\end{figure}

\subsection{Why filtering is the wrong cure}

The broadband overlap is decisive: \emph{no frequency-domain filter on the
gradient can separate the artifact from the signal}, because they share
frequencies. A Helmholtz/Sobolev/Tikhonov smoother does lower the high-to-low
energy ratio, but a faithfulness check --- the cosine of the smoothed gradient
against the true low-mode gradient --- shows it corrupts the very signal it is
meant to preserve. The filter of Section~\ref{sec:filter} is accordingly demoted
from ``the cure'' to ``a mild conditioner usable only atop a structurally smooth
design space.''

The structural cure removes the rough degrees of freedom \emph{by construction}.
We stop treating the boundary nodes as the design and parameterize the hole
radius by a truncated Fourier series
\begin{equation}
  r(\theta) = r_0 + \sum_{m\in\mathcal M}\big(a_m\cos m\theta + b_m\sin m\theta\big),
  \label{eq:fourier-hole}
\end{equation}
with the area held fixed by solving for $r_0$. The noisy nodal gradient is
\emph{contracted} onto the design modes by the exact chain rule,
$\partial C/\partial a_m = \sum_j g_{\mathrm{rad},j}\cos m\theta_j$, which
integrates the per-node noise out. This is a Hilbertian (Riesz) gradient in a
finite-dimensional smooth design space; the adjoint of
Sections~\ref{sec:adjoint-theory}--\ref{sec:five-steps} is untouched --- only the
final contraction $\mat J^\top\vect g$ is new.

\subsection{The clean-DOF budget and its auto-calibration}

A design space helps only if its modes sit \emph{below the artifact floor}. The
admissible mode count is set by the cloud, not by hand. The artifact decorrelates
at the stencil radius $\rho$ (the largest neighbour distance in a stencil), so a
boundary mode is trustworthy only if its half-wavelength spans at least one
stencil radius --- a stencil-Nyquist criterion
$m \le \lfloor \pi r/\rho \rfloor$. The Sobolev length of the residual
conditioner is likewise tied to the cloud, $\ell_{\mathrm{sob}} = \rho/r$. Both
knobs are derived, not tuned:

\begin{center}
\begin{tabular}{lll}
\hline
knob & formula & value ($\mathrm{d}x{=}0.05$, $k{=}35$) \\
\hline
stencil radius $\rho$ & $\max$ neighbour distance & $0.30 = 6\,\mathrm{d}x$ \\
mode cap $m_{\mathrm{cap}}$ & $\lfloor \pi r/\rho \rfloor$ & $2$ \\
Sobolev length & $\rho/r$ & $1.06$ \\
\hline
\end{tabular}
\end{center}

At this resolution the clean-DOF budget is a single mode --- the ellipse mode.
Hand-set richer truncations ($m{=}2{:}4$) are unstable: the optimizer random-walks
the under-resolved $m{=}3,4$ channels up until a boundary wiggle breaks a stencil.
The calibrated cap closes them by construction; unlocking higher modes requires
refinement (smaller $\rho/r$), a resolution/cost trade rather than a tuning knob.

\subsection{Three failure modes, and why they need different cures}
\label{sec:plate-three}

With the gradient exact and the design space clean, the binding obstacle changed
identity. Three distinct failures, conflated in the earlier work, must be kept
apart:

\begin{description}
\item[(A) Descent-direction noise.] The $\partial\mat W/\partial\pts$ roughness
  above. Intrinsic to the discrete gradient --- present on a \emph{fresh} cloud at
  iteration zero --- and cured only by the design-space contraction.
\item[(B) Cloud degradation.] As nodes move, a cloud frozen at an old geometry
  ill-conditions: frozen stencils stretch, frozen adjacency goes stale. A
  stale-cloud artifact.
\item[(C) Objective bias.] A cloud fitted to a \emph{stale} geometry has its
  discrete-compliance optimum displaced from the true one --- on the plate,
  $a_2^\star\approx\pm0.07$ rather than the true circle $a_2=0$, the sign
  bracketed between a frozen interior ($+0.078$) and a Laplace-morphed interior
  ($-0.065$). The shallow optimum (compliance varies $\sim\!2\%$ across the range)
  amplifies the $O(h^p)$ objective error into an $O(0.07)$ design error. Also a
  stale-cloud artifact.
\end{description}

That (C) is a stale-cloud artifact and not intrinsic was settled by a fresh-cloud
test: at the true circle, a cloud freshly generated for that geometry has
$\partial C/\partial a_2 \approx 5\times10^{-8}$, three orders below the driving
gradient --- the circle \emph{is} the discrete optimum of an appropriately
re-initialized cloud. Averaged over independent fresh clouds the residual $m{=}2$
gradient maps to $a_2^\star\approx0.008$, a tenfold debiasing.

\begin{figure}[h]
\centering
\includegraphics[width=0.9\linewidth]{fig_bias.png}
\caption{The objective bias is a \emph{stale-cloud} artifact, not intrinsic. The
discrete-compliance optimum $a_2^\star$ is displaced from the true circle
($a_2{=}0$) and the sign \emph{brackets}: a frozen interior gives $+0.078$, a
Laplace-morphed interior $-0.065$. Re-anchored re-meshing (two-front) lands at
$+0.015$ --- essentially the circle.}
\label{fig:plate-bias}
\end{figure}

A practical corollary closes a tempting shortcut: the old per-node route could
\emph{not} have been rescued by periodic re-initialization alone. Re-meshing cures
(B) and (C) but not (A); on a fresh cloud the per-node gradient is still
Nyquist-dominated, so the boundary roughens immediately. Re-meshing helps the
per-node route only if it \emph{smooths} the boundary --- which is the
parameterization in disguise. (A) and (B,~C) need different cures; the
parameterization is not optional.

\subsection{The two-front architecture}
\label{sec:plate-twofront}

The cures compose into a single loop with two cooperating fronts on the cloud,
the design space handling (A):

\paragraph{Front 1 --- morph (within a re-mesh interval).} The interior is slaved
to the boundary by a smooth, differentiable Laplace extension,
$\Delta\pts_{\mathrm{int}} = \mat M\,\Delta\pts_{\mathrm{bnd}}$ with
$\mat M = -\mat L_{ii}^{-1}\mat L_{ib}$, factorized once per anchor. The adjoint
flows through it by the transpose,
$\hat{\vect g}_{\mathrm{bnd}} = \vect g_{\mathrm{bnd}} + \mat M^\top
\vect g_{\mathrm{int}}$; omitting this term makes the boundary-only gradient
$O(1)$ wrong once the interior moves. Fixed stencils within the interval keep the
objective smooth.

\begin{figure}[h]
\centering
\includegraphics[width=0.6\linewidth]{fig_morph.png}
\caption{Front~1 (morph). As the hole moves from the start ellipse (red) toward
the circle (green), the differentiable Laplace extension carries the interior nodes
inward (displacement arrows, $\times3$), decaying smoothly with distance. The
adjoint flows back through this map by its transpose.}
\label{fig:plate-morph}
\end{figure}

\paragraph{Front 2 --- re-anchored re-mesh (between intervals).} When a quality
indicator trips, a fresh high-quality cloud is generated for the current design
and the morph re-anchored to it (reference reset, operator refactorized). Each
re-mesh reads as a \emph{re-initialization}, not a perturbed continuation --- this
is what removes (B) and (C). Re-meshing is not precluded by the discrete adjoint;
it is precluded only if one differentiates \emph{through} the re-mesh or demands
exact stationarity of a single fixed discrete objective. We do neither --- we
differentiate the smooth morph within an interval and re-anchor between intervals.

Because the discrete compliance \emph{value} jumps at each re-mesh (by an amount
bounded by the discretization accuracy --- noise, not signal), descent is on the
gradient, $\lVert\mat J^\top\vect g\rVert\to0$, never gated on $C$. A
gradient-decaying ``settling'' step was tried and rejected: under a stale morph it
lingers in the distorted regime and corrupts the gradient.

\subsection{Closed-loop cloud quality: which indicators matter}
\label{sec:plate-indicators}

Front~2 is triggered by a registry of cloud-quality indicators, each a cheap
scalar measured every iteration, individually toggleable with its own threshold;
only enabled indicators trigger a re-mesh. On the ellipse$\to$circle deformation:

\begin{center}
\begin{tabular}{lll}
\hline
indicator & measures & verdict \\
\hline
\texttt{morph\_drift} & max interior displacement $/\,\mathrm{d}x$ & active trigger; conservative \\
\texttt{stencil\_growth} & frozen-stencil radius $/$ anchor radius & principled but slack ($\le1.09$) \\
\texttt{min\_gap} & closest interior--hole $/\,\mathrm{d}x$ & dormant \\
\texttt{spacing\_cv} & interior nearest-neighbour CV & dormant \\
\texttt{boundary\_cv} & hole-node spacing CV & dormant \\
\texttt{min\_sep} & global minimum separation $/\,\mathrm{d}x$ & dormant \\
\hline
\end{tabular}
\end{center}

The smooth Laplace morph is gentle: a $0.58\,\mathrm{d}x$ node drift stretches the
worst stencil by only $\sim\!9\%$. The displacement proxy \texttt{morph\_drift}
fires first and is therefore conservative; \texttt{stencil\_growth}, tied directly
to operator validity, is the more principled trigger and would let intervals run
longer. The remaining indicators are insurance for harsher regimes (a growing
hole would wake \texttt{min\_gap}; a strongly non-circular target,
\texttt{boundary\_cv}). \texttt{stencil\_growth} captures ``old neighbours
separate'' but not ``non-neighbours approach'' (adjacency-membership staleness),
which needs a current $k$-NN query and was negligible here.

\subsection{Result}

The two-front loop converges to the circle: $a_2:0.0986\to0.0154$ (true $0$),
$r_0:0.2741\to0.2826$ (true $0.2828$), compliance
$5.18\times10^{-5}\to4.88\times10^{-5}$, with two \texttt{morph\_drift}-triggered
re-meshes over sixty iterations. It lands at the true optimum, not at either
stale-cloud bias ($+0.078$ or $-0.065$): the design space removes the noise, the
re-anchoring removes the bias, and the morph keeps the cloud differentiable in
between.

\begin{figure}[h]
\centering
\includegraphics[width=0.95\linewidth]{fig_twofront.png}
\caption{Two-front result. \emph{Top-left:} hole shapes (start ellipse $\to$ final,
against the target circle). The optimizer drives $a_2\to0$ and the compliance down,
while the cloud-quality indicators (bottom) trigger two re-meshes
(\texttt{morph\_drift}). It converges to the true circle, not to either stale-cloud
bias.}
\label{fig:plate-twofront}
\end{figure}

\section{Differentiation strategy: manual pull-backs and an AD-native API}
\label{sec:diffstrategy}

The pipeline of Sections~\ref{sec:adjoint-theory}--\ref{sec:five-steps} is a discrete
adjoint \emph{orchestrated by hand}: the RBF weights are differentiated automatically
by \texttt{RadialBasisFunctions.jl}, but the chain rule through the assembled system
(the five steps) is written explicitly inside \texttt{shape\_gradient} rather than
traced by AD. This section
explains why, and how the same numerics could be repackaged behind a clean
reverse-mode API without paying the cost that motivated the manual route. It is an
open engineering direction, not yet built.

\subsection{Why not full reverse-mode AD}

The obvious alternative is to write the forward computation in plain Julia and let a
reverse-mode AD engine (Mooncake) differentiate it. We tried this first and
abandoned it: AD \emph{traces every elementary operation}, and the forward solve
contains an UMFPACK factorization and, inside \texttt{\_build\_weights}, one small
RBF saddle solve per stencil. Tracing all of that produced multi-minute LLVM
compile times even at $N=25$ nodes. The manual pipeline sidesteps this by calling
the analytic pull-backs (\texttt{\_pullback\_weights!}, the IFT adjoint solve)
directly --- no AD trace, no compile blow-up.

\subsection{Primitive versus rule}

The cost is avoidable \emph{without} giving up AD, and the distinction that makes
this possible is worth stating plainly because it is easy to conflate:
\begin{itemize}
  \item A \emph{rule} (reverse rule, \texttt{rrule}) is the analytic derivative of
        one function: given $y=f(x)$ it returns $y$ and a \emph{pull-back}
        $\bar x = \mathrm{pb}(\bar y)$. It is the mathematics.
  \item A function is a \emph{primitive} when AD is told to \emph{use the rule
        instead of tracing inside it}. One makes a function a primitive \emph{by}
        giving it a rule.
\end{itemize}
So ``primitive'' and ``rule'' are not alternatives --- a primitive is a function
equipped with a rule. The leverage: declaring the two expensive functions
(\texttt{build\_weights} and \texttt{solve}) primitives makes AD stop at their
boundary and call our \emph{existing} analytic pull-backs, while it traces only the
cheap glue (geometry parameterization, block assembly, scalar loss).

\subsection{The schematic}

The forward computation reads like ordinary code; the annotations mark which nodes
AD traces and which are primitives backed by an analytic rule:

\begin{footnotesize}
\begin{verbatim}
FORWARD (reads like normal code):

  theta --geometry--> pts --build_weights--> W
        --assemble--> (A,b) --solve--> u --loss--> C

  geometry, assemble, loss : AD-traced (cheap glue)
  build_weights            : PRIMITIVE, rule = _pullback_weights!  (dW/dx)
  solve                    : PRIMITIVE, rule = IFT adjoint solve

BACKWARD (Mooncake runs this once the primitives are declared):

  C_bar --loss--> (u_bar, b_bar)
        --[solve rule: one adjoint solve  A^-T (dC/du)]--> (A_bar, b_bar)
        --assemble--> W_bar
        --[build_weights rule: _pullback_weights! = dW/dx]--> pts_bar
        --geometry--> theta_bar
\end{verbatim}
\end{footnotesize}

Only the primitive nodes carry expensive numerics, and those are exactly the
hand-written pull-backs already validated to $10^{-7}$ against finite differences.
Everything else is small, so AD compiles quickly.

\subsection{What changes, concretely}

\begin{center}
\begin{tabular}{lll}
\hline
piece & today (manual) & AD-native version \\
\hline
\texttt{build\_weights} & \texttt{\_pullback\_weights!} by hand & primitive $+$ \texttt{rrule!!} wrapping it \\
\texttt{solve} & IFT adjoint in \texttt{shape\_gradient} & primitive $+$ \texttt{rrule!!} (forward $+$ adjoint solve) \\
assembly, loss, geometry & hand-coded external terms & plain Julia, AD-traced \\
call site & \texttt{shape\_gradient(pts, model, N, \ldots)} & \texttt{gradient(loss, theta)} \\
\hline
\end{tabular}
\end{center}

Two recurring pains vanish in the AD-native form. First, the design parameterization
composes automatically: \texttt{gradient(theta -> loss(geometry(theta)))} chains
$\partial C/\partial\pts \cdot \partial\pts/\partial\theta$ --- the parameterization
Jacobian that Section~\ref{sec:plate} showed is essential for convergence --- with no
extra code. Second, the two-term $\partial\vect b/\partial\pts$ subtlety
(Section~\ref{sec:step5-load}: the explicit load-sensitivity term that must be added
by hand in the manual API) disappears, because writing $C=\vect b^\top\vect u$ lets
AD trace the loss's $\vect b$-dependence itself.

\subsection{Status --- an open engineering point}

This is a \emph{repackaging of validated numerics, not new mathematics}, but it is
not yet done and carries real engineering risk. Registering a function as a Mooncake
primitive incurs a one-time \texttt{rrule!!} build cost (already the dominant cost in
the present shape-opt loop), which must be amortized through the prebuilt sysimage;
and the analytic pull-backs mutate cached buffers, so wrapping them under Mooncake's
calling convention needs care. The recommended next step is a small timed prototype
--- wrap \texttt{build\_weights} and \texttt{solve} as primitives, check
\texttt{gradient(loss, theta)} against the manual \texttt{shape\_gradient} and against
finite differences, and measure the warm-up --- before the framework extraction
commits to an API. The pay-off, if the warm-up is acceptable, is the industry-facing
simplicity of a single \texttt{gradient} call in place of the present
multi-argument entry point.

\section{Generalization}
\label{sec:generalization}

The five-step pipeline is structured so that each step is
PDE-agnostic except where the PDE explicitly enters. We summarize
what changes when the PDE does.

\subsection{Scalar elliptic PDEs}

For scalar problems (heat equation, Poisson, electrostatics,
magnetostatics) the system is $N \times N$, not $2N \times 2N$.
Steps 1, 2, 4 are essentially unchanged. Step 3 simplifies: the
coefficient extraction in~\eqref{eq:interior-extract} reduces from
the four-coefficient family to a single per-operator constant.
The Neumann row family becomes $\partial u/\partial \vect n = g$
or $\alpha u + \beta \partial u/\partial\vect n = g$ (Robin), with
a one-coefficient analogue of~\eqref{eq:neumann-row}.

\subsection{Nonlinear PDEs (Newton iteration)}

For nonlinear $\vect F(\vect u, \pts) = \vect 0$ solved by Newton,
the adjoint is built at the \emph{converged} Jacobian:
\begin{equation}
  \mat J(\vect u^\star)^\top \vect\eta
  \;=\; \big(\partial\Loss/\partial\vect u\big)^\top
  \quad \text{at the converged iterate } \vect u^\star.
\end{equation}
Steps 1, 2, 4 carry over identically. Step 3 changes: the
coefficient extraction now has solution-dependent coefficients
(the Jacobian entries involve $\vect u^\star$ itself), and the
$\partial\mat F/\partial\pts$ contracts both with state and with
the operator weight matrices. The detailed treatment for
incompressible Navier--Stokes is sketched in
\texttt{plan\_manual\_adjoint.md}.

\subsection{Time-dependent PDEs}

Transient problems require a backward-in-time sweep of the
adjoint:
\begin{equation}
  -\frac{\dd\vect\eta}{\dd t}
  \;+\; \mat J^\top \vect\eta
  \;=\; \big(\partial\Loss/\partial\vect u\big)^\top
  \quad \text{from } t=T \text{ back to } t=0.
\end{equation}
The architectural cost is checkpointing (memory for storing
forward states at intermediate times) but the per-step structure
is the same as the steady case.

\subsection{Multi-physics coupling}

Coupling several physics on the same point cloud --- the
fundamental advantage of the meshless approach for this project
--- is block-structural in the algebra: the joint $\mat A$ has
diagonal blocks for each physics and off-diagonal blocks for
couplings, all built from the same RBF stencils. The Step 3
coefficient extraction generalizes to per-block coefficient
families. The Step 4 backward through the weight matrices is
unchanged. Strongly-coupled problems with a coupled forward
Newton solve fall under the nonlinear-PDE case above.

\subsection{3D}

The five-step pipeline is dimension-agnostic in its
\emph{structure}; only the per-operator content changes. In 3D:
\begin{itemize}[leftmargin=2em]
  \item the interior weight matrix family grows to
        $\{\partial_{xx}, \partial_{yy}, \partial_{zz},
        \partial_{xy}, \partial_{yz}, \partial_{xz}\}$;
  \item the Neumann row family in~\eqref{eq:neumann-row}
        becomes a $3 \times 3$ analogue (with 3 normal
        components and 3 stress components);
  \item the polyline normal in~\eqref{eq:chord-normal} is
        replaced by a triangle-based vertex-normal-average
        (sparse on the three vertices of each STL face), with
        the analogous closed-form Jacobian;
  \item the boundary-loop Helmholtz filter is replaced by a
        Laplacian on the discrete surface mesh.
\end{itemize}
None of these changes requires new mathematical content. The
content is the same; the bookkeeping is more.

\section{Position relative to the state of the art}
\label{sec:sota}

\subsection{Reference points}

The mature reference points for shape optimization in
contemporary engineering practice are FEM-based commercial codes
(Tosca, OptiStruct, NX Topology, ANSYS Mechanical Topology),
density-based topology optimization via SIMP/RAMP
(Bendsøe \& Sigmund, codified in the 88-line MATLAB code of
Andreassen et al.\ 2011), CAD-parameterized shape optimization
(ESP, Vorpal, OpenFOAM's continuous-adjoint solvers), and the
level-set / phase-field methods of Wang et al.\ 2003 onward. The
workhorse optimizer is the Method of Moving Asymptotes
(Svanberg), with IPOPT, L-BFGS-B, and SQP rounding out the
toolbox. Helmholtz-PDE sensitivity filtering for SIMP was
introduced by Lazarov \& Sigmund (Int.\ J.\ Numer.\ Methods Eng.,
2011) --- essentially the filter we use here, applied to a 1D
boundary loop rather than to the design density field.

\subsection{Honest comparison}

The present work is best described as \emph{methodologically
sound on a niche discretization, not yet operationally
competitive}. Table~\ref{tab:sota} compares dimensions where the
gap is most evident.

\begin{table}[h]
\centering
\caption{Honest comparison of the present pipeline against the
state of the art.}
\label{tab:sota}
\renewcommand{\arraystretch}{1.3}
\begin{tabular}{p{0.30\textwidth}p{0.30\textwidth}p{0.32\textwidth}}
\toprule
Dimension & State of the art & This work \\
\midrule
PDE discretization & FEM at $10^6$ DOFs; AMR; multi-physics
  & RBF-FD at $N \sim 45$ in the example; single physics
    (lin.\ elasticity); 2D \\
Adjoint correctness check & FD at small scale, otherwise
checkpointed AD
  & Central order-5 FD vs analytic at every config; $\sim 10^{-7}$ \\
Boundary normal differentiability & implicit through CAD/FEM mesh
  & Explicit chord formula with sparse closed-form Jacobian
    (Section~\ref{sec:l2}) \\
Sensitivity filter & Helmholtz PDE on the design field
(Lazarov--Sigmund 2011) or radial average; arc-length / surface
  & Helmholtz PDE on the 1D boundary loop, uniform segments
    (Section~\ref{sec:filter}) \\
Optimizer & MMA, IPOPT, L-BFGS-B, SLP, SQP
  & Steepest descent with Armijo backtracking \\
Constraint handling & Augmented Lagrangian, projected gradient,
SQP
  & Quadratic penalty with hand-tuned $\rho_{\text{pen}}$ \\
Geometry kernel & CAD/NURBS/splines or STL-vertex direct,
auto-meshing
  & Boundary points are the design (no kernel yet); STL absent \\
Mesh quality under deformation & Re-mesh / ALE / FFD
  & Not needed (meshless) \\
Multi-material / topology change & SIMP density field, level set
  & Out of scope of this report; meshless naturally permits it \\
Scale (DOFs) & $10^4$ to $10^7$, 3D
  & $\sim 10^2$, 2D \\
\bottomrule
\end{tabular}
\end{table}

\subsection{Where the work might claim novelty}

A defensible claim would be along the lines of: \emph{a complete
manual-adjoint shape-optimization pipeline for RBF-FD linear
elasticity, with closed-form analytic sensitivities at the
weight-matrix and normal levels and a Helmholtz boundary-loop
sensitivity filter, validated to $10^{-7}$ against finite
differences and structured for clean extension to multi-physics
on the same point cloud}. RBF-FD has been used extensively for
forward PDE solving (Bayona, Flyer, Fornberg and collaborators),
but its application to shape optimization with analytic discrete
adjoints is sparse in the literature; most existing RBF-FD
shape-opt work relies on black-box AD or continuous-adjoint
formulations that ignore discretization error. The deliberate
separation of geometry (\texttt{WhatsThePoint.jl}) from PDE
machinery (\texttt{Macchiato.jl}) is also unusual and
well-motivated for meshless contexts.

This is a methodologically clean starting point in an
under-explored corner of the field. It is not state-of-the-art
performance, scale, or feature coverage; we will return to the
gap in Section~\ref{sec:openwork}.

\section{Open work}
\label{sec:openwork}

\subsection{Within the elasticity pipeline}

\paragraph{Stronger optimizer.}
Steepest descent with Armijo converges slowly on anisotropic
Hessians. The natural next step is L-BFGS-B for the
box-constrained version of the problem and MMA for the full
multi-constraint version. Both are wrapper-level work --- the
adjoint pipeline is unchanged.

\paragraph{Augmented Lagrangian for the area constraint.}
The current quadratic penalty produces a sawtooth in
$\norm{\nabla\Loss}$ as $V(\pts)$ oscillates around $V_0$ across
line-search trials. An augmented Lagrangian, with multiplier
updated between outer iterations, eliminates this artefact.

\paragraph{Filter tuning --- moot.}
The adaptive-$r$ and arc-length-Laplacian refinements of the boundary filter are no
longer pursued: Section~\ref{sec:plate} supersedes gradient filtering with the
design-space contraction.

\paragraph{Grid resolution study.}
$N = 45$ is too coarse to express a clean Michell taper. A scale
test at $N \gtrsim 100$ would clarify whether the regularized
optimum converges to a clean tapered beam as the grid refines.

\subsection{Within the cantilever example}

\paragraph{Cantilever retired.}
The cantilever has served its purpose (adjoint validation, regularization study) and
is retired; the plate-with-hole (Section~\ref{sec:plate}) is the current case. The
illustrative numbers in Section~\ref{sec:cantilever} carry a known tip-shrinkage sign
subtlety (masked by an earlier corner-freeze convention) and are not re-validated ---
the qualitative contrasts they show are what we use them for.

\paragraph{Reference textbook problem --- addressed.}
Beyond FD-vs-AD consistency, the discretization itself is now checked against an
analytic solution: the Kirsch plate-with-hole stress concentration
($K_t = 2.966$ at $0.62\%$ L2 stress error).

\subsection{Toward engineering geometries}

\paragraph{Geometry kernel.}
The design is now a smooth, low-dimensional parameterization $\vect l$
(Section~\ref{sec:plate}), not the boundary points directly. For
real STL-driven geometries the differentiable chain
$\vect l \to \mathcal G_{\text{STL}} \to \pts_b$ is needed, with
$\vect l$ either (i) STL vertices direct, (ii) spline / NURBS
control points, or (iii) FFD lattice control points. This is the
single gating item between the present R\&D state and engineering
applicability.

\paragraph{STL-side primitives in \texttt{WhatsThePoint.jl}.}
The cross-repo separation places STL operations, point cloud
generation, and triangle-based vertex normals on the WTP side.
Once the geometry kernel choice is made, the WTP work is: STL
vertex perturbation, triangle-normal derivation with the analogue
of~\eqref{eq:chord-jac}, STL-to-pointcloud sampling with
Jacobian, optional spline-to-STL retessellation.

\paragraph{Boundary-to-interior deformation --- done.}
When boundary points move, interior points must follow to maintain cloud quality.
This is now Front~1 (Section~\ref{sec:plate-twofront}): a differentiable Laplace
morph, with the adjoint carried by its transpose. The remaining open item is its
\emph{efficient} re-factorization between re-meshes at large $N$.

\paragraph{Scale-up.}
Forward and adjoint solves at $N \sim 10^4$--$10^5$ require
iterative or preconditioned linear solvers tuned for RBF-FD
elasticity. No off-the-shelf preconditioner exists; building one
is a research-grade subproblem.

\subsection{Beyond elasticity}

\paragraph{Other physics, same pipeline.}
The five-step pipeline as derived in Section~\ref{sec:five-steps}
is reusable for any PDE that fits the
$\mat A(\pts)\vect u = \vect b(\pts)$ template, with adapted
coefficient extraction in Step 3. Heat conduction, electrostatics,
magnetostatics, and incompressible Stokes are direct extensions;
incompressible Navier--Stokes requires the Newton-converged
Jacobian variant (Section~\ref{sec:generalization}).

\paragraph{Multi-physics shape optimization.}
The architectural advantage of meshless is that the same point
cloud supports all of these physics with the same weight matrix
family. Coupling is block-structural; the adjoint of the coupled
system is built by the same Step 2 solve, with the joint right-
hand side from each physics' contribution to the loss. The
strategic motivation for the framework, articulated in
\texttt{Vision.md}, is precisely this multi-physics tractability.

\section{Conclusion}

We have presented the complete discrete adjoint pipeline
implemented in \texttt{Macchiato.jl} for RBF-FD shape optimization
of plane-stress linear elasticity. Starting from the continuous
shape-optimization problem, we derived the five-step decomposition
in which the sensitivity calculation reduces to: a forward solve
that caches the RBF stencils; a single adjoint solve; an analytic
back-projection from $\Delta\mat A = -\vect\eta\vect u^\top$ onto
per-operator sensitivities $\Delta\mat W^{\mathcal L}$ via the
elasticity coefficient structure (interior) and the Neumann
coefficient family (boundary); a pure-Julia backward pass through
the RBF stencils via \texttt{\_pullback\_weights!}; and a closed
set of explicit external contributions covering the area
constraint, the shape-dependent traction values (Phase D
Level 1), and the differentiable polyline normals (Phase D
Level 2). The whole chain is validated to $\sim 10^{-7}$ relative
error against fifth-order central finite differences.

The cantilever motivated a sensitivity filter on the boundary loop,
but the plate-with-hole study of Section~\ref{sec:plate} showed that
frequency filtering is the wrong cure: the boundary gradient noise is
the geometry sensitivity of the RBF-FD weights,
$\partial\mat W/\partial\pts$, which is broadband and overlaps the
signal, so no gradient filter separates them faithfully. The cure is
structural --- a clean-DOF Fourier design space (calibrated to the
stencil-resolved mode budget) that removes the rough degrees of freedom
by construction, an actively re-anchored re-mesh that removes the
stale-cloud objective bias, and a differentiable Laplace morph that
keeps the cloud smooth in between. On the plate the resulting two-front
optimizer converges to the true circular optimum, landing at neither
stale-cloud bias; the filter survives only as a mild conditioner atop
the smooth design space.

The framework is structured for clean extension to other physics
on the same point cloud, to three dimensions with vertex-normal
gradients on STL triangles, and to multi-physics shape
optimization --- the strategic target of the broader project.
Each ingredient is in place; what remains is engineering
extension rather than mathematical invention. The path from this
working document to a publishable paper runs through the geometry
kernel decision and a representative engineering test case;
those are the next milestones, and the present pipeline is the
foundation they will rest on.

\end{document}
